\section{Critical laws after Gaussian regularization}\label{sec:critical}

Suppose that the conclusion of Theorem~\ref{thm:main} is false.
This section constructs, from a sequence of counterexamples, smooth
probabilities $\eta$ satisfying a logarithmic Sobolev inequality and
vector fields $w$ that nearly satisfy $A_\eta w=w$. There are two
possible constructions. An attained nonlinear entropy quotient gives
a score $w$; equality of the logarithmic Sobolev and Poincar\'e
constants gives the gradient of a gap eigenfunction. We shall retain
both alternatives throughout.

\subsection{Normalization and Gaussian regularization}

\begin{lemma}[Normalized counterexamples]\label{lem:critical-normalization}
Failure of Theorem~\ref{thm:main} produces a sequence of centered,
compactly supported, full dimensional log-concave probabilities $\mu$
such that
\begin{equation}\label{eq:critical-normalization}
 R(\mu)=1,\qquad
 \sup_{|\theta|=1}\E_\mu e^{\langle X,\theta\rangle^2/a^2}
 \le2,\qquad a\longrightarrow0.
\end{equation}
For an absolute constant $C$, they also satisfy
\begin{equation}\label{eq:critical-linear-mgf}
 \log\E_\mu e^{v\cdot X}\le Ca^2|v|^2,
 \qquad \Cov_\mu(X)\le Ca^2\Id.
\end{equation}
\end{lemma}

\begin{proof}
Choose laws and compact smooth tests for which the entropy-to-energy
ratio divided by the squared marginal parameter tends to infinity.
The tests may be chosen with positive energy: a smooth test of zero
energy is constant on the connected interior of the convex support,
and its entropy vanishes. Restriction of the law to a sufficiently
large ball preserves each test quotient arbitrarily closely. Multiply
the restricted density by $e^{-\varepsilon|x|^2/2}$ and normalize,
with $\varepsilon>0$ sufficiently small to preserve that quotient
again. Both operations preserve log-concavity, their combined density
relative to the original law can be bounded by $2$, and
Lemma~\ref{lem:curvature} gives a finite logarithmic Sobolev
constant for the resulting compact law.

A bounded density increase changes the exponential-square expectation
by at most that factor. If $\E e^{Z^2/A^2}\le4$, Jensen gives
$\E e^{Z^2/(2A^2)}\le2$. Centering costs only another absolute
factor: $\E Z^2\le A^2\log4$, while
$(Z-\E Z)^2\le2Z^2+2(\E Z)^2$; increase the denominator first
to make the resulting exponential expectation bounded by an absolute
constant, and once more by Jensen to make it at most $2$.
These arguments apply separately to every unit linear form with
the same constants. Translation preserves the test quotient. Thus
the centered compact laws still have $R/A^2\to\infty$ with
$0<R<\infty$. Dilation by $R^{-1/2}$ gives
Equation~\eqref{eq:critical-normalization}.

For completeness, if a centered scalar $Z$ obeys
$\E e^{Z^2/a^2}\le2$, then
$\E|Z|^k\le2a^k\Gamma(k/2+1)$, by integrating its tail.
Expansion of $e^{tZ}$, with the linear term absent, therefore gives
$\log\E e^{tZ}\le Ca^2t^2$ for $a|t|\le c$.
For $a|t|\ge c$, use
$tZ\le Z^2/a^2+a^2t^2/4$ and absorb $\log2$ into
$Ca^2t^2$. Apply this to $Z=X\cdot\theta$. The second moment
bound follows either from the exponential-square bound or by
differentiating the centered moment generating function at zero.
\end{proof}

We suppress the counterexample index, but all limits in this section
refer to this sequence. Let $G$ be a fresh standard Gaussian and set
\[
 \mu_s=\Law(X+\sqrt sG),\quad R_s=R(\mu_s),\quad
 b_s(y)=\E[X\mid X+\sqrt sG=y],\quad V_s=-\log\mu_s.
\]
Here and below $\mu_s$ also denotes its strictly positive density.
Gaussian differentiation and posterior covariance give
\begin{equation}\label{eq:critical-posterior}
 DV_s(y)=\frac{y-b_s(y)}s,\qquad
 M_s(y)=Db_s(y)=\frac{\Cov(X\mid y)}s,\qquad
 K_s(y)=D^2V_s(y)=\frac{\Id-M_s(y)}s.
\end{equation}
The posterior of $X$ has a convex potential plus $|x-y|^2/(2s)$.
Lemma~\ref{lem:curvature}, applied to its linear forms, therefore
implies
\begin{equation}\label{eq:critical-curvature}
 0\le M_s\le\Id,\qquad 0\le K_s\le s^{-1}\Id.
\end{equation}
All these functions are smooth. Compact support of $X$ also gives
$|b_s|\le\sup_{\supp\mu}|x|$.

The monotonicity below follows from the heat-flow contraction of
Klartag and Putterman \cite[Theorem~1.2 and Section~3]{KlartagPutterman2023},
which uses the transport construction of Kim and Milman
\cite{KimMilman2012}. We recall the argument for our regularized laws.

\begin{lemma}[Bounds for the regularized constants]
\label{lem:critical-flow-bounds}
The function $s\mapsto R_s$ is nondecreasing and $1$-Lipschitz on
$(0,\infty)$, and
\begin{equation}\label{eq:critical-basic-R}
 1\le R_s\le1+s.
\end{equation}
\end{lemma}

\begin{proof}
The heat equation for $\mu_s$ is the continuity equation with
velocity $DV_s/2$. On a compact positive-time interval this velocity
is globally Lipschitz and grows at most linearly, so its smooth flow
is a diffeomorphism carrying $\mu_s$ to $\mu_t$. Its Jacobian
satisfies $\dot J=K_tJ/2$. By positivity of $K_t$, the flow
expands distances and its inverse is $1$-Lipschitz when $t\ge s$.
Transport a test by that inverse. Its entropy is unchanged and
its Dirichlet energy cannot increase, proving $R_t\ge R_s$.
Equation~\eqref{eq:calculus-convolution} gives
$R_{s+h}\le R_s+h$ and $R_s\le1+s$. For every compact smooth
test, its entropy and energy converge to those under $\mu$ as
$s\downarrow0$. Taking the supremum of its quotient gives
$\liminf_{s\downarrow0}R_s\ge R(\mu)=1$. Monotonicity now
proves the remaining bound.
\end{proof}

\subsection{Attainment above the Gaussian and Poincar\'e thresholds}

On compact manifolds, Rothaus obtains nonconstant positive solutions
of the logarithmic Euler equation when the sharp logarithmic Sobolev
constant exceeds the Poincar\'e constant \cite[p.~358]{Rothaus1986}.
The proof below must also control loss of entropy in the Gaussian tails.
The next two lemmas establish this control and attainment. Their constants
may depend on a single fixed law and its dimension. Uniform estimates
along the counterexample sequence begin with the translation argument.

\begin{lemma}[Compactness and the Gaussian entropy defect]
\label{lem:critical-defect}
Fix $s>0$ and a compactly supported probability $\mu$ on $\R^n$,
and put $\nu=\mu*\gamma_s$. Its form domain embeds compactly
into $L^2(\nu)$. For every $\varepsilon>0$ there is a finite
$C_\varepsilon$, depending on this fixed law, such that
\begin{equation}\label{eq:critical-defective}
 \Ent_\nu(u^2)\le2(s+\varepsilon)\|Du\|_2^2
                  +C_\varepsilon\|u\|_2^2.
\end{equation}
If $u_j\to u$ in $L^2(\nu)$, $\|u_j\|_2=1$, and
$\|Du_j\|_2^2\to E_*<\infty$, then
\begin{equation}\label{eq:critical-entropy-defect}
 \limsup_j\Ent_\nu(u_j^2)
 \le\Ent_\nu(u^2)+2s(E_*-\|Du\|_2^2).
\end{equation}
\end{lemma}

\begin{proof}
Write $r=d\nu/d\gamma_s$. The Gaussian kernel formula gives
$D\log r=b_s/s$, which is bounded, and hence
$|\log r(y)|\le C(1+|y|)$. For $v=u\sqrt r$,
\[
 \Ent_\nu(u^2)=\Ent_{\gamma_s}(v^2)
                         -\int v^2\log r\,d\gamma_s,
 \qquad Dv=\sqrt r\,(Du+\tfrac12uD\log r).
\]
Gaussian integration by parts followed by Cauchy--Schwarz gives
\[
 \int |y|^2v^2d\gamma_s
 =sn\|v\|_2^2+2s\int v\,y\cdot Dv\,d\gamma_s
 \le2sn\|v\|_2^2+4s^2\|Dv\|_2^2.
\]
Initially this calculation is for compact tests, and cutoffs extend
it by lower semicontinuity. The linear-growth entropy correction
is consequently bounded by
$\delta\|Dv\|_2^2+C_\delta\|v\|_2^2$ for every $\delta>0$.
Gaussian logarithmic Sobolev and the displayed derivative formula,
with another arbitrarily small Cauchy--Schwarz parameter, prove
Equation~\eqref{eq:critical-defective}. The same moment estimate
shows tightness of the $L^2(\nu)$ mass of any bounded form
sequence. On a fixed ball, the smooth density is bounded above
and away from zero, so the ordinary local Rellich theorem gives
$L^2$ precompactness. A diagonal extraction and tightness prove
compactness globally.

We give the tail argument for the sharp coefficient $2s$ in
Equation~\eqref{eq:critical-entropy-defect}. Put
$v_{j,L}=(|u_j|-L)_+$ and $m_{j,L}=\|v_{j,L}\|_2^2$.
Strong $L^2$ convergence implies
$\lim_{L\to\infty}\sup_jm_{j,L}=0$. The energies of
$v_{j,L}$ and $|u_j|\wedge L$ sum to $\|Du_j\|_2^2$.
Lower semicontinuity for the clipped functions therefore gives
\begin{equation}\label{eq:critical-tail-energy}
 \lim_{L\to\infty}\limsup_j\|Dv_{j,L}\|_2^2
 \le E_*-\|Du\|_2^2.
\end{equation}
For a nonnegative $v$ of squared norm $m\le1$, the negative
part of $v^2\log v^2$ has integral at most $C\sqrt m$
(use $-z\log z\le C\sqrt z$ on $0\le z\le1$).
Since $m\log m\le0$, Equation~\eqref{eq:critical-defective}
thus gives
\[
 \int v^2\log_+v^2d\nu
 \le2(s+\varepsilon)\|Dv\|_2^2
          +C_\varepsilon m+C\sqrt m.
\]
If $|u_j|>ML$, where $M>1$ is fixed, then
$|u_j|\le c_M v_{j,L}$ with $c_M=M/(M-1)$. Consequently
\[
 u_j^2\log_+u_j^2
 \le c_M^2v_{j,L}^2\log_+v_{j,L}^2
                +c_M^2\log(c_M^2)v_{j,L}^2.
\]
The integrals of bounded truncations of $z^2\log z^2$ pass to
the limit by strong $L^2$ convergence; its negative part does as
well, being bounded and continuous. The last two bounds and
Equation~\eqref{eq:critical-tail-energy}, with limits in the order
$j$, $L$, $M$, $\varepsilon$, show that the lost positive
entropy is at most $2s(E_*-\|Du\|_2^2)$. The normalization
terms are zero because $\|u\|_2=\|u_j\|_2=1$.
\end{proof}

The threshold $s$ in the entropy-defect estimate controls escape into the
Gaussian tails.  The Poincar\'e constant controls maximizing sequences
that converge to constants.  The next lemma proves attainment when the
logarithmic Sobolev constant exceeds both thresholds.

\begin{lemma}[Existence and positivity of an optimizer]\label{lem:optimizer}
Let $\nu=\mu*\gamma_s$, where $s>0$ and $\mu$ is compactly
supported. Suppose its logarithmic Sobolev and Poincar\'e constants
$R=R(\nu)$ and $P=R^P(\nu)$ satisfy
$\infty>R>\max(P,s)$. Then there is a smooth strictly positive
nonconstant $f$ with
\begin{equation}\label{eq:critical-optimizer}
 \int f^2d\nu=1,\qquad
 \Ent_\nu(f^2)=2R\int|Df|^2d\nu,\qquad
 A_\nu f=\lambda f\log f,\quad\lambda=R^{-1}.
\end{equation}
The last equation is initially distributional and then classical.
\end{lemma}

\begin{proof}
Take nonnegative unit-norm maximizing tests $f_j$. Replacing a test
by its absolute value preserves entropy and energy, so this choice
loses nothing. Equation~\eqref{eq:critical-defective}, with
$s+\varepsilon<R$, bounds their energies. We first exclude
energies tending to zero; this is the only place the strict
inequality $R>P$ is required.

Suppose $\varepsilon_j^2=\|Df_j\|_2^2\to0$, and write
$c_j=\E_\nu f_j$ and $h_j=(f_j-c_j)/\varepsilon_j$.
Poincar\'e gives $\|h_j\|_2^2\le P$, while
$\|Dh_j\|_2^2=1$ and
$1-c_j^2=\varepsilon_j^2\|h_j\|_2^2$.
Thus $c_j\to1$. By compactness, after extraction,
$h_j\to h$ in $L^2$, where
\[
 \E h=0,\qquad \|Dh\|_2^2\le1,\qquad
 \|h\|_2^2\le P\|Dh\|_2^2.
\]
For $c>0$, $\varepsilon>0$ and $z\ge-c/\varepsilon$ define
\[
 Q_{c,\varepsilon}(z)=\varepsilon^{-2}
 \left[(c+\varepsilon z)^2\log\frac{(c+\varepsilon z)^2}{c^2}
                  -(c+\varepsilon z)^2+c^2\right].
\]
Normalization and $-\log c_j^2-1+c_j^2\ge0$ give
\begin{equation}\label{eq:critical-small-energy-start}
 \varepsilon_j^{-2}\Ent(f_j^2)
 \le\int Q_{c_j,\varepsilon_j}(h_j)\,d\nu.
\end{equation}
Here is the scalar estimate needed to control unbounded $h_j$.
For $1/2\le c\le1$ and $0<\varepsilon\le1/2$,
\begin{equation}\label{eq:critical-scalar-Q}
 0\le Q_{c,\varepsilon}(z)
 \le(2+2\log2)z^2+z^2\log_+z^2,
 \qquad Q_{c,\varepsilon}(z)\longrightarrow2z^2
\end{equation}
uniformly on bounded $z$-intervals as $c\to1$ and
$\varepsilon\to0$. To verify the bound, put $t=\varepsilon z/c$.
Then $Q=z^2B(t)$, where Taylor's integral formula gives
\[
 B(t)=2+4\int_0^1(1-r)\log(1+rt)\,dr,
 \qquad -1\le t<\infty.
\]
For $t\le0$, $B(t)\le2$. For $t\ge0$, use
$t\le z$ to obtain
$B(t)\le2+2\log(1+z)\le2+2\log2+\log_+z^2$.
Continuity at $t=0$ proves the compact convergence.

The tail proof in Lemma~\ref{lem:critical-defect} did not use
unit norm until its last sentence. In particular, for any bounded
form sequence $u_j\to u$ in $L^2$ with energies tending to $E_*$,
it proves
\begin{equation}\label{eq:critical-general-tail}
 \lim_{K\to\infty}\limsup_j
 \int_{|u_j|>K}u_j^2\log_+u_j^2d\nu
 \le2s(E_*-\|Du\|_2^2).
\end{equation}
Apply the compact convergence in Equation~\eqref{eq:critical-scalar-Q}
on $|h_j|\le K$ and its upper bound on the complement. Strong
$L^2$ convergence removes the quadratic tail as $K\to\infty$;
Equation~\eqref{eq:critical-general-tail} controls the entropy tail.
Equation~\eqref{eq:critical-small-energy-start} consequently gives
\[
 \limsup_j\varepsilon_j^{-2}\Ent(f_j^2)
 \le2\|h\|_2^2+2s(1-\|Dh\|_2^2)
 \le2\max(P,s)<2R.
\]
This contradicts the maximizing property.

We may therefore extract $f_j\to f$ strongly in $L^2$ with
$\|Df_j\|_2^2\to E_*>0$. Put $E=\|Df\|_2^2$. The entropy
defect and the sharp inequality give
\[
 2RE_*\le\Ent(f^2)+2s(E_*-E)
         \le2RE+2s(E_*-E).
\]
Since $R>s$ and $E\le E_*$, equality forces $E=E_*>0$.
Thus $f$ is a nonnegative, nonconstant optimizer.

The homogeneous functional
$\Ent(g^2)-2R\|Dg\|_2^2$ is nonpositive and is zero at $f$.
Its derivative along every compact smooth variation $\zeta$ is
\[
 4\int f\log f\,\zeta\,d\nu
       -4R\int Df\cdot D\zeta\,d\nu=0.
\]
The entropy normalization term cancels the additional linear term.
Differentiation is legitimate also at zeros because
$z^2\log z^2$ has derivative zero at zero. This proves the Euler
equation in Equation~\eqref{eq:critical-optimizer}.

We have obtained a nonnegative optimizer with positive energy.  To use
its logarithmic derivative, we still need local regularity and strict
positivity.  We spell out the weak-solution regularity step, using the
interior Newtonian-potential $L^p$ and Schauder estimates; see
\cite[Lecture~24, p.~2; Lecture~10, Proposition~1]{Viaclovsky2004}.
Write $\nu=e^{-V}$ and set $v=e^{-V/2}f\ge0$.  The function $V$ is smooth,
and direct differentiation of the weak equation gives
\[
 -\Delta v=\lambda v\log v+
 \left(\frac{\lambda V}{2}-\frac{|DV|^2}{4}
                         +\frac{\Delta V}{2}\right)v.
\]
The conjugating factor and its reciprocal are bounded on each fixed
ball.  Since $|z\log z|\le C_\delta(1+z^{1+\delta})$ for $z\ge0$,
local $L^p$ integrability of $v$ puts the right side in
$L^q_{\rm loc}$, where $q=p/(1+\delta)>1$.  Localize that right side
on a larger ball and take its Newtonian potential with the sign solving
$-\Delta u=F$.  The cited $L^q$ estimate gives $u\in W^{2,q}_{\rm loc}$;
the difference $v-u$ is distributionally harmonic and hence smooth by
Weyl's lemma \cite[Lecture~3, Theorem~2]{Viaclovsky2004}.
This proves the required $W^{2,q}$ gain directly for the weak solution.

For $n\ge2$, start with $p=2$ and choose
$0<\delta<\min(1,2/n)$.  Before boundedness is reached, Sobolev
embedding improves reciprocal integrability from $1/p$ to
$(1+\delta)/p-2/n$, a decrease greater than $1/n$ because $p\ge2$.
At a critical exponent use any sufficiently large finite integrability
exponent and continue once more.  After finitely many steps $q>n/2$,
so $v$, and thus $f$, is locally bounded and H\"older continuous.
In dimension one the same conclusion follows by twice integrating the
local Poisson equation.  The function $z\log z$ is locally H\"older
of every exponent below one on $[0,\infty)$; reducing the exponent
if necessary, the Schauder estimate therefore gives
$f\in C^{2,\alpha}_{\rm loc}$ for some $\alpha>0$.
All constants in this paragraph concern one fixed law and dimension;
none enters a uniform estimate along the counterexample sequence.

If $f(x_0)=0$, choose a ball $B$ about $x_0$ with
$0\le f\le b<1$ on its closure. Start the $\nu$-diffusion at
$x_0$ and stop at its exit time $\tau$ from $B$. For
$Z_t=f(Y_{t\wedge\tau})$, $m(t)=\E Z_t$, and
$g(z)=-\lambda z\log z$, the stopped It\^o formula gives, for almost
every $t$,
\[
 m'(t)=\E[\1_{t<\tau}g(Z_t)]\le\E g(Z_t)\le g(m(t)),
 \qquad m(0)=0.
\]
Indeed $g$ is nonnegative and concave on $[0,b]$; the first
inequality adds its nonnegative stopped boundary contribution.
Also $m'\ge0$. If $m(t_1)>0$, let $t_a$ be the first time at
which $m=a$, for $0<a<m(t_1)$. Integration gives
\[
 \int_a^{m(t_1)}\frac{dz}{-\lambda z\log z}
 \le t_1-t_a\le t_1,
\]
contradicting divergence as $a\downarrow0$. Hence $m=0$.
The killed nondegenerate diffusion has positive probability of
reaching each open subset of $B$ while staying in $B$: one can
follow a smooth interior path by requiring Brownian motion to
remain in a small tube about its corresponding driving path,
and use local Lipschitz continuity of the drift. Brownian tube
probabilities are positive by its finite increment densities and
bridge continuity. Consequently $f$ vanishes throughout $B$.
Its zero set is both open and closed, impossible in connected
$\R^n$ for a unit-norm function. Thus $f>0$, and smooth elliptic
bootstrap makes $f$ smooth everywhere.
\end{proof}

An optimizer supplies a second probability with its own logarithmic
Sobolev inequality. The next identity explains why no convexity
assumption on this second probability is needed.

\begin{samepage}
\begin{lemma}[Change of reference at an optimizer]
\label{lem:critical-change-reference}
Under Lemma~\ref{lem:optimizer}, suppose in addition that
$\nu=e^{-V}$ is log-concave with bounded $K=D^2V$.
Define
\[
 \eta=f^2\nu,\qquad h=\log f^2,\qquad w=Dh,\qquad
 I=\E_\eta|w|^2.
\]
Then $I=2\lambda H(\eta\mid\nu)>0$, $R(\eta)\le1/\lambda$,
each component of $w$ belongs to $\Dom(A_\eta)$, and
\begin{equation}\label{eq:critical-score-equation}
 A_\nu h=\lambda h+\tfrac12|w|^2,
 \qquad A_\eta w=(\lambda\Id-K)w.
\end{equation}
Moreover,
\begin{equation}\label{eq:critical-curvature-mean}
 \E_\eta\langle w,Kw\rangle\le\lambda|\E_\eta w|^2.
\end{equation}
\end{lemma}
\end{samepage}

\begin{proof}
Since $w=2Df/f$, $I=4\|Df\|_{L^2(\nu)}^2$; optimizer
equality gives $I=2\lambda H(\eta\mid\nu)>0$.
The chain rule applied to its Euler equation gives the first
identity in Equation~\eqref{eq:critical-score-equation}. Commuting
one derivative with $A_\nu$ gives
\[
 A_\nu w=D(A_\nu h)-Kw
       =\lambda w+(Dw)w-Kw.
\]
Since $A_\eta=A_\nu-w\cdot D$ and $Dw$ is symmetric, the
second identity follows distributionally. Its right side is in
$L^2(\eta)$, since $K$ is bounded and $w\in L^2(\eta)$.
Lemma~\ref{lem:dirichlet-core} therefore puts $w$ in the
operator domain; no unproved integrability of higher derivatives
has been used.

Let $p=g^2\eta$ have unit mass, where $g$ is compactly supported
and smooth. Direct expansion and compactly supported integration
by parts give
\begin{align*}
 \Ical(p\mid\nu)
 &=\Ical(p\mid\eta)+\int g^2|w|^2d\eta
                    +2\int D(g^2)\cdot w\,d\eta\\
 &=\Ical(p\mid\eta)+\int g^2\{|w|^2+2A_\eta h\}\,d\eta\\
 &=\Ical(p\mid\eta)+2\lambda\int h\,dp.
\end{align*}
These formulas may be read in square-root form at zeros of $g$.
On the other hand
$H(p\mid\nu)=H(p\mid\eta)+\int h\,dp$. Thus
\begin{equation}\label{eq:critical-exact-deficit}
 \Ical(p\mid\nu)-2\lambda H(p\mid\nu)
 =\Ical(p\mid\eta)-2\lambda H(p\mid\eta).
\end{equation}
Apply the inequality for $\nu$ to the compact smooth function
$gf$. Its deficit is nonnegative, so the right side is
nonnegative. Homogeneity and the Sobolev approximation of
Section~\ref{sec:calculus} prove $R(\eta)\le1/\lambda$.
Finally the domain identity gives
\[
 \E_\eta|Dw|^2=\lambda I-\E_\eta\langle w,Kw\rangle,
\]
while Poincar\'e for $\eta$ gives the lower bound
$\lambda(I-|\E_\eta w|^2)$. Subtract to obtain
Equation~\eqref{eq:critical-curvature-mean}.
\end{proof}

\subsection{Uniform estimates in the two alternatives}

\begin{lemma}[Translation gain]\label{lem:translation-gain}
Let $\mu$ satisfy Equation~\eqref{eq:critical-normalization}, with
$a\le1$. Fix a compact interval $J\subset(0,1/2)$. There is a
constant $C_J$, independent of the dimension and of the support
radius of $\mu$, with the following property. For $s\in J$, put
$\lambda=R_s^{-1}$. If $\sigma$ is any probability with
\[
 D_0=H(\sigma\mid\mu_s)<\infty,\qquad
 F_0=\Ical(\sigma\mid\mu_s)<\infty,\qquad p=\E_\sigma Y,
\]
then
\begin{equation}\label{eq:critical-translation-gain}
 F_0-2\lambda D_0
 \ge(s^{-2}-\lambda/s)|p|^2
                   -C_J\sqrt a\,(F_0+D_0).
\end{equation}
\end{lemma}

\begin{proof}
The centered linear moment generating functions under $\mu_s$
obey
\[
 \log\E_{\mu_s}e^{v\cdot Y}\le C_J|v|^2,
 \qquad
 \log\E_{\mu_s}e^{v\cdot b_s(Y)}\le Ca^2|v|^2.
\]
The first follows by independent Gaussian addition and
Equation~\eqref{eq:critical-linear-mgf}; the second follows by
conditional Jensen from that equation. Entropy duality, optimized
over the scalar multiple of each unit vector, gives
\begin{equation}\label{eq:critical-entropic-means}
 |p|\le C_J\sqrt{D_0},\qquad
 |\E_\sigma b_s(Y)|\le Ca\sqrt{D_0}.
\end{equation}
All expectations exist. More explicitly, compact support of $X$
makes $\E_{\mu_s}e^{c|Y|^2}<\infty$ for some $c>0$;
entropy duality then implies $\E_\sigma|Y|^2<\infty$.
The latter fact is used only for integrability, not for a uniform
bound in dimension.

Let $\sigma_u$ be the translate of $\sigma$ by $-up$,
$0\le u\le1$, and let $D_u=H(\sigma_u\mid\mu_s)$.
Change of variables in entropy and the Hessian bound give
\begin{align*}
 D_u
 &=D_0+\E_\sigma[V_s(Y-up)-V_s(Y)]\\
 &\le D_0-u|p|^2/s+u\,p\cdot\E_\sigma b_s(Y)/s
                                +u^2|p|^2/(2s)
 \le C_JD_0.
\end{align*}
Thus Equation~\eqref{eq:critical-entropic-means}, applied to
$\sigma_u$, implies
$|\E_\sigma b_s(Y-up)|\le C_Ja\sqrt{D_0}$ uniformly in $u$.
Differentiating the exact entropy expression, and then integrating,
gives
\begin{equation}\label{eq:critical-translated-entropy}
 D_1=D_0-\frac{|p|^2}{2s}+O_J(aD_0).
\end{equation}
Indeed its derivative is
$-(1-u)|p|^2/s+p\cdot\E_\sigma b_s(Y-up)/s$.

Write $r(Y)=D\log(d\sigma/d\mu_s)(Y)$ in the weak
square-root sense. Its mean is
\begin{equation}\label{eq:critical-score-mean}
 \E_\sigma r=\E_\sigma DV_s=(p-\E_\sigma b_s)/s.
\end{equation}
To justify this for finite Fisher information, the Lebesgue density
$q$ of $\sigma$ has weak derivative
$Dq=q(r-DV_s)\in L^1$: $r$ is square integrable and $DV_s$
has linear growth. Integrating against expanding cutoffs gives
$\int Dq=0$ and hence Equation~\eqref{eq:critical-score-mean}.

In the moving coordinates $Y$, the score of $\sigma_1$ relative
to $\mu_s$ is
\[
 r(Y)-p/s+d(Y)/s,\qquad d(Y)=b_s(Y)-b_s(Y-p).
\]
Since $0\le Db_s\le\Id$, integration along the segment gives
\[
 |d(Y)|^2
 =\left|\int_0^1M_s(Y-up)p\,du\right|^2
 \le\int_0^1|M_s(Y-up)p|^2du
 \le p\cdot d(Y).
\]
The mean estimates at the two endpoints therefore imply
\begin{equation}\label{eq:critical-posterior-translation-error}
 \E_\sigma|d|^2\le C_JaD_0.
\end{equation}
Expanding the score square and using
Equation~\eqref{eq:critical-score-mean},
\[
 \E_\sigma|r-p/s|^2
 =F_0-|p|^2/s^2+2p\cdot\E_\sigma b_s/s^2.
\]
The remaining cross term is bounded by
$C_J\sqrt{aD_0}\,(\sqrt{F_0}+\sqrt{D_0})$, and the
$d$-square term by $C_JaD_0$. Consequently the translated
information is finite and satisfies
\begin{equation}\label{eq:critical-translated-information}
 F_1=F_0-|p|^2/s^2+O_J(\sqrt a\,(F_0+D_0)).
\end{equation}
These calculations require only the weak derivatives already
justified, so they apply also when the density of $\sigma$ vanishes.
Apply logarithmic Sobolev to $\sigma_1$, namely
$F_1\ge2\lambda D_1$, and combine
Equations~\eqref{eq:critical-translated-entropy} and
\eqref{eq:critical-translated-information}. Since
$2/3\le\lambda\le1$ on the time interval in question, this
proves Equation~\eqref{eq:critical-translation-gain}.
\end{proof}

We can now obtain estimates that are uniform along the counterexample
sequence. Write $R_s^P$ for the optimal Poincar\'e constant of
$\mu_s$. By linearization $R_s^P\le R_s$. Since $R_s\ge1>s$
on $(0,1/2)$, the strict alternative $R_s>R_s^P$ meets all
the hypotheses of Lemma~\ref{lem:optimizer}.

\begin{lemma}[Critical fields in both alternatives]
\label{lem:critical-two-alternatives}
Fix $J\subset(0,1/2)$ compact and $s\in J$, and put
$\lambda=1/R_s$. There are a smooth positive probability $\eta$
and a smooth vector field $w$, with
$0<I:=\E_\eta|w|^2<\infty$, such that
\begin{equation}\label{eq:critical-common-identities}
 R(\eta)\le\lambda^{-1},\qquad
 w\in\Dom(A_\eta)^n,\qquad
 A_\eta w=(\lambda\Id-K_s)w.
\end{equation}
They can be chosen as follows.
\begin{enumerate}[label=\textup{(\roman*)}]
\item If $R_s>R_s^P$, take the optimizer $f$ of
Lemma~\ref{lem:optimizer} and set
\[
 \eta=f^2\mu_s,\qquad w=D\log f^2,
 \qquad I=2\lambda H(\eta\mid\mu_s).
\]
\item If $R_s=R_s^P$, there is a centered nonconstant gap
eigenfunction $\varphi$ satisfying
$A_{\mu_s}\varphi=\lambda\varphi$. Take
\[
 \eta=\mu_s,\qquad w=D\varphi,
 \qquad I=\lambda\E_{\mu_s}\varphi^2.
\]
\end{enumerate}
In either alternative,
\begin{equation}\label{eq:critical-uniform-smallness}
 |\E_\eta w|^2+\E_\eta|K_sw|^2\le C_J\sqrt a\,I,
 \qquad
 \E_\eta\langle w,K_sw\rangle\le\lambda|\E_\eta w|^2.
\end{equation}
In alternative \textup{(ii)} the factor $\sqrt a$ can be
replaced by $a^2$. Coordinates belong to $\Dom(A_\eta)$ in
both alternatives.
\end{lemma}

\begin{proof}
In the first alternative the identities and the last inequality
are Lemma~\ref{lem:critical-change-reference}. Apply
Lemma~\ref{lem:translation-gain} to $\sigma=\eta$, where
$F_0=I$ and $D_0=I/(2\lambda)$, so the deficit vanishes.
The coefficient $s^{-2}-\lambda/s$ is bounded below by a
positive constant on $J$. Thus, with $p=\E_\eta Y$,
\[
 |p|^2\le C_J\sqrt a\,I,
 \qquad |\E_\eta b_s|\le C_Ja\sqrt I.
\]
Equation~\eqref{eq:critical-score-mean} applied to $\eta$ now
gives $|\E_\eta w|^2\le C_J\sqrt a\,I$.
Since $K_s^2\le s^{-1}K_s$, the curvature-mean inequality
gives the asserted estimate for $K_sw$.

In the second alternative, compact form embedding supplies a
minimizer of the Rayleigh quotient among mean-zero functions of
unit norm. Its positive eigenvalue is $1/R_s^P=\lambda$.
The nullspace consists of constants, since zero weak gradient on
$\R^n$ implies constancy. Compactness therefore makes this a
positive gap, and interior elliptic regularity makes its
eigenfunction $\varphi$ smooth. Commuting a derivative gives
$A_{\mu_s}D\varphi=(\lambda\Id-K_s)D\varphi$.
Its right side is in $L^2$, so
Lemma~\ref{lem:dirichlet-core} gives the operator-domain assertion.
Taking the energy pairing and then Poincar\'e gives, exactly as
before,
\[
 \E\langle w,K_sw\rangle\le\lambda|\E w|^2.
\]
Coordinates belong to the operator domain because both $Y$ and
$A_{\mu_s}Y=(Y-b_s)/s$ belong to $L^2(\mu_s)$. Integration
by parts, or self-adjointness with these coordinate tests, yields
\[
 \E w=\lambda\E[\varphi Y]
      =\frac{\E[\varphi Y]-\E[\varphi b_s]}s,
 \qquad
 (1-\lambda s)\E[\varphi Y]=\E[\varphi b_s].
\]
Conditional Jensen gives
\[
 \E_{\mu_s}b_sb_s^*\le\E_\mu XX^*\le Ca^2\Id.
\]
For every unit vector $v$, Cauchy--Schwarz thus bounds
$|\E[\varphi\,v\cdot b_s]|$ by $Ca\|\varphi\|_2$.
Taking the supremum over $v$ introduces no dimension factor.
Since $1-\lambda s\ge1/2$ and
$I=\lambda\|\varphi\|_2^2$, it follows that
$|\E w|\le C_Ja\sqrt I$. The bound for $K_sw$ follows
from $K_s^2\le s^{-1}K_s$ once again.

It remains to justify the coordinate assertion for the first
alternative. Finite $H(\eta\mid\mu_s)$ and the Gaussian
tails of $\mu_s$ imply $\E_\eta|Y|^2<\infty$, by the
entropy argument in the proof of Lemma~\ref{lem:translation-gain}.
The distributional identity
\begin{equation}\label{eq:critical-coordinate-generator}
 A_\eta Y=(Y-b_s)/s-w
\end{equation}
has an $L^2(\eta)$ right side, because $b_s$ is bounded for
the fixed law and $w\in L^2(\eta)$. The maximal-domain
assertion of Lemma~\ref{lem:dirichlet-core} applies.
\end{proof}

\begin{proposition}[Limit of the regularized constants]
\label{prop:regularized-constants}
For the sequence in Lemma~\ref{lem:critical-normalization},
\begin{equation}\label{eq:critical-R-limit}
 R_s\longrightarrow1\qquad\text{for every }0<s<1/2.
\end{equation}
At every differentiability point of $R_s$, the fields of
Lemma~\ref{lem:critical-two-alternatives} satisfy the contact
identity
\begin{equation}\label{eq:critical-contact-derivative}
 R_s'=R_s\,\frac{\E_\eta\langle w,K_sw\rangle}{I}.
\end{equation}
In particular $0\le R_s'\le C_J\sqrt a$ almost everywhere
on each fixed compact $J\subset(0,1/2)$.
\end{proposition}

\begin{proof}
Fix $s$, and let $\Phi_{s,t}$ be the heat continuity flow from
Lemma~\ref{lem:critical-flow-bounds}, for $t$ on either side
of $s$. Transport a form-domain test $g$ by
$g_t=g\circ\Phi_{s,t}^{-1}$. Its distribution under $\mu_t$
is fixed, so both its entropy and its variance are constant.
The Jacobian equation $\partial_tD\Phi_{s,t}=K_tD\Phi_{s,t}/2$
gives
\begin{equation}\label{eq:critical-transport-energy}
 \left.\frac{d}{dt}\right|_{t=s}
 \int|Dg_t|^2d\mu_t
 =-\int\langle Dg,K_sDg\rangle d\mu_s.
\end{equation}
For compact smooth $g$ this follows by differentiating the
inverse-transpose Jacobian in the gradient. For a form-domain
$g$, the change-of-variables formula still holds by Sobolev
approximation. The flow derivatives and their time derivatives
are uniformly bounded on a compact time interval, so the
derivative of the integrand is dominated by $C|Dg|^2$.
Dominated convergence proves
Equation~\eqref{eq:critical-transport-energy} without a smoothness
or compact-support restriction on $g$.

In the strict alternative take $g=f$ and set
$Q(t)=\Ent_{\mu_t}(f_t^2)/(2\int|Df_t|^2d\mu_t)$.
The denominator is positive near $s$, and $Q(t)\le R_t$ for
both signs of $t-s$, with $Q(s)=R_s$. At a differentiability
point of $R_s$, this two-sided contact implies $Q'(s)=R_s'$.
Equation~\eqref{eq:critical-transport-energy} and
$w=2Df/f$ give Equation~\eqref{eq:critical-contact-derivative}.
In the gap alternative take $g=\varphi$ and instead use
$Q(t)=\Var_{\mu_t}(\varphi_t)/\int|D\varphi_t|^2d\mu_t$.
Now $Q(t)\le R_t^P\le R_t$ and $Q(s)=R_s^P=R_s$.
The same contact argument gives the identical derivative formula,
with $w=D\varphi$.

Equation~\eqref{eq:critical-uniform-smallness} bounds this
derivative by $C_J\sqrt a$ in either alternative, with constants
independent of which alternative holds at each time. Lipschitz
continuity makes $R_s$ absolutely continuous. For fixed
$0<\delta<s<1/2$, integration on $[\delta,s]$ gives
\[
 1\le R_s\le R_\delta+C_{[\delta,s]}\sqrt a\,(s-\delta)
             \le1+\delta+C_{[\delta,s]}\sqrt a\,(s-\delta).
\]
First take the counterexample-sequence limit and then let
$\delta\downarrow0$. This proves
Equation~\eqref{eq:critical-R-limit}.
\end{proof}

We finish by fixing the Gaussian time and collecting precisely the
objects used in the following sections. The fixed time will not be
sent to zero.

\begin{proposition}[The critical sequence]\label{prop:critical-sequence}
If Theorem~\ref{thm:main} fails, there is a sequence of dimensions
and centered compactly supported log-concave laws $\mu$ satisfying
Equation~\eqref{eq:critical-normalization}. Fix $s=1/8$, form
$\mu_s,b_s,M_s,K_s$ as in Equation~\eqref{eq:critical-posterior},
and put $\lambda=R_s^{-1}$. After passage to a subsequence there
are smooth positive probabilities $\eta$ and smooth vector fields
$w$ such that, with $A=A_\eta$ and $I=\E_\eta|w|^2$,
\begin{gather}
 \lambda\longrightarrow1,\qquad R(\eta)\le\lambda^{-1},
 \qquad 0<I<\infty,\qquad w\in\Dom(A)^n,
 \label{eq:critical-sequence-data}\\
 A w=(\lambda\Id-K_s)w,\qquad
 \|(A-1)w\|_{L^2(\eta)}+|\E_\eta w|=o(\sqrt I),
 \label{eq:critical-near-eigenfunction}\\
 \E_\eta|K_sw|^2=o(I),\qquad
 \E_\eta|Dw|^2=(1+o(1))I.
 \label{eq:critical-sequence-energy}
\end{gather}
The spectral gap of $A$ is at least $\lambda$. Each coordinate
function belongs to $\Dom(A)$. Exactly one of the following
constructions can be used throughout the subsequence:
\begin{enumerate}[label=\textup{(\roman*)}]
\item $R_s>R_s^P$, $\eta=f^2\mu_s$ for the positive unit-norm
optimizer, $w=D\log f^2$, and
$I=2\lambda H(\eta\mid\mu_s)$. The exact deficit identity in
Equation~\eqref{eq:critical-exact-deficit} holds for compact square-root
tests relative to $\eta$, and
$A Y=(Y-b_s)/s-w$.
\item $R_s=R_s^P$, $\eta=\mu_s$,
$A\varphi=\lambda\varphi$, $\E_\eta\varphi=0$,
$w=D\varphi$, and $I=\lambda\E_\eta\varphi^2$.
One may rescale $\varphi$ so that $I\to0$.
\end{enumerate}
All $o$-terms refer to the counterexample sequence with fixed
finite parameters. In particular no relation between dimension,
$I$, and the rate of convergence is asserted or needed here.
\end{proposition}

\begin{proof}
Choose an infinite subsequence with the same alternative at
$s=1/8$. Lemma~\ref{lem:critical-two-alternatives} gives all
structural identities and the curvature and mean bounds.
Proposition~\ref{prop:regularized-constants} gives $\lambda\to1$.
Consequently
\[
 (A-1)w=(\lambda-1)w-K_sw=o_{L^2}(\sqrt I).
\]
The energy pairing yields
$\E|Dw|^2=\lambda I-\E\langle w,K_sw\rangle=(1+o(1))I$.
Poincar\'e with constant $\lambda^{-1}$ is the asserted spectral
gap. The coordinate and branch-specific identities have already
been proved. In the gap alternative all estimates are homogeneous
in $\varphi$, so multiplying it by positive scalars can enforce
$I\to0$ without changing any normalized assertion.
\end{proof}
